\section{Konzeption der Modellarchitektur}
\label{ch:konzeption_modellarchitektur}

Aufbauend auf der systematischen Transformation des unstrukturierten Rohdatensatzes in einen maschinenlesbaren Zustandsraum erfolgt in diesem Kapitel die informationstechnische Konzeption der neuronalen Modellarchitektur. 
Das primäre methodische Ziel besteht in der Konstruktion einer Policy $\pi_\theta$, welche in der Lage ist, auf Basis historischer Zeitreihen einen kontinuierlichen Aktionsraum -- spezifisch die Leistungs- und Bremshebelstellung -- prädiktiv abzubilden. 

\subsection{Auswahlder Architektur}
\label{sec:auswahl_architektur}

Um eine fundierte Entscheidung für die finale Modellstruktur zu treffen, wurde ein initiales experimentelles Setup aufgesetzt. 
In dieser Evaluationsphase wurden ein \acrlongshort{MLP}, ein \acrlongshort{LSTM} sowie eine \acrlongshort{GRU} vergleichend trainiert und analysiert.

\subsubsection{Experimentelles Setup}
Die initiale Trainingsphase wurde bewusst ohne Regularisierungsmaßnahmen wie Dropout oder Weight Decay (L2-Regularisierung) durchgeführt. 
Dieses Vorgehen stellt sicher, dass das native Lernverhalten sowie die prinzipielle Eignung der Architekturtypen für die Verarbeitung fahrphysikalischer Zeitreihen isoliert betrachtet werden können. 
Für eine effiziente Evaluation wurde das Training auf 20\,\% des Gesamtdatensatzes limitiert und über 10 Epochen mit einer Batch-Size von 256 sowie einer Learning Rate von 0.001 ausgeführt. 
Das Eingangsfenster umfasste 100 Zeitschritte ($H=100$), was physikalisch 20 Sekunden an historischen Daten entspricht, um einen Vorhersagehorizont von 50 Zeitschritten (10 Sekunden) zu prädizieren.

Die Auswertung des Feedforward-Netzwerkes verdeutlicht die theoretischen Limitationen dieser Architektur. 
Mit 110.849 trainierbaren Parametern wies das \gls{MLP} die massivste Überdimensionierung im Testfeld auf und erreichte konsequenterweise den geringsten Trainingsfehler (siehe Abbildung \ref{fig:training_loss_comparison}). 
Auf den ungesehenen Evaluierungsdaten versagte das Modell jedoch: 
Der Validierungs-\gls{RMSE} lag bei 0.2297 (\gls{MAE} = 0.1476), und das Bestimmtheitsmaß auf dem Test-Set fiel auf ein $R^2$ von lediglich 0.0467 bei einer Richtungs-Accuracy von 47,52\,\% und einer Toleranz-Accuracy von 70,57\,\%. 

Aus maschineller Lernperspektive handelt es sich hierbei um ein klassisches Phänomen der extremen Überanpassung (Overfitting). 
Wie in den theoretischen Grundlagen hergeleitet, besitzen Feedforward-Netze keinen internen Zustand (Hidden State) und verarbeiten Eingaben rein statisch. 
Sie sind nicht in der Lage, zeitliche Kausalitäten -- wie die immense physikalische Trägheit eines Zuges -- zu erfassen. 
Das \gls{MLP} nutzt seine enorme Parameteranzahl folglich ausschließlich zur reinen Memorization (Auswendiglernen) der Trainingsdaten. 
Sobald das Modell auf ungesehene Evaluierungsdaten trifft, fehlt ihm das temporale Verständnis, was zu einer Explosion des Validierungs-Losses führt.

Im direkten Kontrast dazu demonstrierten die rekurrenten Architekturen eine deutliche Überlegenheit. 
Das \gls{LSTM} (19.009 Parameter) erzielte ein Test-$R^2$ von 0.1030 bei einem \gls{RMSE} von 0.1550 und einem \gls{MAE} von 0.0898. 
Die Richtungs-Accuracy stieg auf 60,95\,\% (Toleranz-Accuracy: 73,55\,\%). 
Die \gls{GRU} (14.273 Parameter) lieferte ähnliche, minimal schwächere Metriken mit einem Test-$R^2$ von 0.0854 (\gls{RMSE} = 0.1565, \gls{MAE} = 0.0864) und einer Richtungs-Accuracy von 60,85\,\%.

\begin{figure}[htbp]
    \centering
    \includegraphics[width=\textwidth]{loss_comparison.png}
    \caption[Trainings- und Validierungsverluste]{Vergleich der Trainings- und Validierungsverluste der evaluierten Basis-Architekturen zur Visualisierung des architekturabhängigen Lernverhaltens.}
    \label{fig:training_loss_comparison}
\end{figure}

Diese Ergebnisse belegen die Theorie: 
\gls{LSTM} und \gls{GRU} glänzen, da sie exakt für derartige hochgradig sequenzielle Aufgaben entworfen wurden. 
Durch ihre rekurrente Architektur und die spezialisierten Gating-Mechanismen (wie den Cell State beim \gls{LSTM}) können sie historische Abhängigkeiten über das gesamte 20-sekündige Eingangsfenster aggregieren. 
Das Modell begreift mathematisch, dass ein aktuelles Bremsmanöver kausal mit infrastrukturellen oder physikalischen Signalen aus der Vergangenheit zusammenhängt.

Ein wissenschaftlich interessantes Phänomen zeigte sich zudem bei den Trainingszeiten. 
Aus informationstheoretischer Sicht sollte die \gls{GRU} recheneffizienter agieren, da sie durch die Einsparung von Gates mit lediglich 14.273 Parametern deutlich schlanker ist als das \gls{LSTM} (19.009 Parameter). 
Dennoch beanspruchte das Training der \gls{GRU} 654,8 Sekunden, während das \gls{LSTM} bereits nach 399,7 Sekunden konvergierte.

Diese Laufzeit-Anomalie lässt sich durch die zugrundeliegende Software-Architektur erklären. 
In etablierten Deep-Learning-Frameworks wie PyTorch wurden die Implementierungen von \gls{LSTM}-Zellen über Jahre hinweg auf Hardware- und C++-Ebene extrem stark optimiert. 
Die internen Matrixmultiplikationen der \gls{LSTM}-Zellen lassen sich oftmals effizienter parallelisieren als die strikt sequenziellen Update-Schritte der \gls{GRU}. 
Diese hochgradige Software-Optimierung überwiegt in der Praxis den theoretischen Rechenvorteil der geringeren Parameterzahl bei Weitem.

\subsubsection{Fazit der Architekturauswahl}

Unter Berücksichtigung der fundamentalen Anforderungen an die Generalisierungsfähigkeit fällt die finale Entscheidung -- trotz des minimal stärkeren Test-$R^2$ des \glsentryshort{LSTM}s im initialen Setup -- auf die \gls{GRU}-Architektur. 
Aus maschineller Lernperspektive bietet die \gls{GRU} durch ihre reduzierte Parameterzahl einen entscheidenden strukturellen Vorteil: Eine geringere Modellkapazität verringert nativ die Anfälligkeit für Overfitting und zwingt das Netzwerk dazu, robustere, abstrahierte Repräsentationen der Fahrphysik zu erlernen. 
Der im Vorfeld identifizierte Laufzeitnachteil der \gls{GRU}, welcher auf PyTorch-spezifische Implementierungsdetails zurückzuführen ist, ist für den weiteren Projektverlauf zu vernachlässigen, da das finale und rechenintensive Modelltraining ohnehin auf einer deutlich leistungsstärkeren Hardware-Infrastruktur durchgeführt wird. 
Da jedoch auch die \gls{GRU} als unregularisiertes Basis-Modell nicht vollständig immun gegen das Auswendiglernen von Trajektorien ist, ist die Definition und Implementierung fundierter Regularisierungsstrategien für die Konstruktion der finalen Policy zwingend erforderlich.

\subsection{Dimensionierung des Netzwerks}
\label{subsec:dimensionierung_netzwerk}

Die Dimensionierung eines neuronalen Netzes unterliegt einem fundamentalen informationstheoretischen Dilemma: 
Eine zu geringe Modellkapazität bedingt Underfitting, wodurch die komplexe Zugdynamik nicht ausreichend abgebildet werden kann. 
Eine zu große Architektur neigt hingegen -- insbesondere bei verrauschten realen Sensordaten -- zu starkem Overfitting, da das Modell den Trainingsdatensatz auswendig lernt (Memorization) und seine Generalisierungsfähigkeit einbüßt. 
Um das optimale Gleichgewicht empirisch zu ermitteln, wurde ein Brute-Force-Grid-Search-Verfahren implementiert. 
Hierbei wurden 39 permutationsfähige Architekturen des ausgewählten \gls{GRU}-Modells mit einer Tiefe von ein bis drei Hidden Layern und Breiten von 64, 96 sowie 128 Neuronen unter identischen Bedingungen (10 Epochen, Batch-Size 128, Learning Rate 0.001) trainiert. 
Als finale Entscheidungsmetrik diente der \acrlongshort{MSE} auf dem isolierten Test-Datensatz.

Die empirische Evaluation der Single-Layer-Architekturen offenbarte das schlechteste Leistungsniveau. 
Dies belegt, dass ein einzelner rekurrenter Layer nicht über die notwendige Repräsentationskapazität verfügt, um die hochdimensionalen Abstraktionen zwischen Streckentopologie, Kinematik und resultierender Hebelstellung über das definierte Zeitfenster ($H=30$) zu erlernen. 
Eine Erhöhung der Tiefe auf drei Hidden Layer (mit bis zu 251.265 Parametern) verschlechterte die Prädiktionsgüte jedoch massiv. 
Diese tiefen Modelle unterliegen starkem Overfitting, da sie aufgrund der enormen Parameterzahl streckenspezifische Eigenheiten und Sensorrauschen memorieren, anstatt die kausale Fahrphysik zu generalisieren.

Das globale Optimum des Experiments verortet sich in den 2-Layer-Architekturen, welche im Durchschnitt den geringsten Test-Fehler (0.0205) aufweisen. 
Als beste Einzelkonfiguration kristallisierte sich ein Netzwerk mit 96 Neuronen im ersten und 64 Neuronen im zweiten Hidden Layer heraus (Test-\gls{MSE}: 0.0197 bei 61.697 Parametern). 
Das finale Modell für das \gls{BC} wird folglich als 2-Layer-\gls{GRU} mit dieser asymmetrischen Trichterstruktur (Funnel-Design) dimensioniert. 
Dieses Layout ist konzeptionell äußerst vorteilhaft: Der breitere erste Layer extrahiert komplexe, hochfrequente zeitliche Kausalitäten aus dem acht-dimensionalen Input-Strom. 
Der darauffolgende, schmalere Layer agiert als ,,Information Bottleneck''. 
Er zwingt das Modell, die extrahierten Informationen auf die essenziellen physikalischen Treiber zu komprimieren, bevor die finale Vorhersageschicht den kontinuierlichen Stellwert generiert. 
Dieses Design bietet das ideale Verhältnis aus Repräsentationskapazität und inhärenter Regularisierung.

\subsection{Loss-Funktion und Optimizer}
\label{subsec:loss_und_optimizer}

Die finale Prädiktionsgüte der neu ausgewählten 2-Layer-\gls{GRU}-Architektur (96-64 Neuronen) wird im hochdimensionalen Parameterraum maßgeblich von der mathematischen Fehlerbewertung (Verlustfunktion) sowie der algorithmischen Optimierungsstrategie determiniert. 
Um die optimale Konfiguration für diese tiefe Netzwerkstruktur zu ermitteln, wurde eine erneute systematische Grid-Search-Evaluation auf dem Test-Datensatz durchgeführt. 
Hierbei wurden etablierte Regressions-Losses (\gls{MSE}, \gls{MAE}, Huber, Log-Cosh, Smooth L1) in Kombination mit klassischen und adaptiven Optimizern (\gls{SGD}, \gls{Adam}, \gls{AdamW}, \gls{RMSprop}, \gls{Adagrad}) getestet. 
Das Training erfolgte analog zur initialen Architekturauswahl rein auf Basis der unregularisierten \gls{GRU} über 10 Epochen und auf 20\,\% des Gesamtdatensatzes, um die Performance-Unterschiede der Algorithmen isoliert und datengetrieben bewerten zu können.

Die empirischen Ergebnisse dieser Evaluation sind in Form einer Heatmap in Abbildung \ref{fig:loss_optimizer_heatmap} visualisiert. 
Die spaltenweise Analyse der Verlustfunktionen zeigt, dass der Huber-Loss im aggregierten Durchschnitt über alle Optimizer hinweg die höchste Stabilität und den besten mittleren Test-\gls{MSE} (0.0247) aufweist. 
Dennoch liefert der klassische Standard-\gls{MSE} in Kombination mit dem \gls{Adam}-Optimizer den absolut geringsten Einzel-Fehlerwert des gesamten Experiments (Test-\gls{MSE} von 0.0236). 
Klassische Fehlerfunktionen zeigen sich ansonsten anfälliger für die Wahl des Optimizers, wie die deutlichen Schwankungen in der Matrix belegen.

\begin{figure}[htbp]
    \centering
    \includegraphics[width=\textwidth]{heatmap_test_mse.png}
    \caption{Heatmap der Test-\gls{MSE}-Werte für die systematische Grid-Search-Evaluation unterschiedlicher Loss-Funktionen und Optimizer-Strategien auf Basis der unregularisierten 2-Layer-\gls{GRU}-Architektur.}
    \label{fig:loss_optimizer_heatmap}
\end{figure}

Die zeilenweise Auswertung der Optimierungsalgorithmen offenbart ein noch prägnanteres Leistungsgefälle. Das klassische Stochastische Gradientenabstieg-Verfahren (\gls{SGD}) versagt in diesem Setup unabhängig von der gewählten Loss-Funktion komplett und weist die signifikant schlechteste Generalisierung auf (dunkelviolette Zeile mit Werten von bis zu 0.0279 beim \gls{MSE}). 
Adaptive Algorithmen dominieren das Feld sichtbar und stabilisieren das Training auf einem durchgehend höheren Niveau. 
\gls{Adagrad} besticht dabei durch eine enorme Konstanz über alle Fehlerfunktionen hinweg und erzielt den besten mittleren Test-\gls{MSE} der Optimizer-Familien (0.0241). 
\gls{RMSprop} zeigt mit \gls{MAE} und Log-Cosh (jeweils unter 0.0240) ebenfalls extrem starke Resultate. 
Dennoch erreicht die \gls{Adam}-Variante in der Spitze das absolute globale Optimum.

Basierend auf dieser direkten Auswertung der Heatmap wird das finale \gls{BC}-Modell mit dem \gls{Adam}-Optimizer und dem \gls{MSE}-Loss kompiliert. 
Diese Konfiguration stellt empirisch belegt die maximale Vorhersagegüte für die nun tiefer dimensionierte Netzwerkstruktur sicher.


